\section*{Abstract}

British Columbia (BC) sets an allowable annual cut (AAC), a regulatory ceiling, while fire-resilient
species composition emerges from enterprise regeneration decisions. We ask whether the AAC is
feasible under deciduous-forward post-clearcut stocking standards and what those standards cost in
value and tail risk. We formulate a bilevel evaluation: the regulator chooses a post-clearcut stocking standard, and the
enterprise chooses an even-flow schedule under the AAC. We apply it to the Prince George TSA and search for even-flow schedules under four standards: the historical clearcut and D26, D55, and D84 (26, 55,
and 84 percent deciduous targets). Every period is at or above the 2025 AAC
(\SI{6.917231}{\mega\cubic\meter\per\year}); the final even-flow reference period reaches
\SI{7.28}{\mega\cubic\meter\per\year} (\SIrange{5.2}{5.3}{\percent} above the ceiling), and the mean
is \SIrange{6.935}{6.938}{\mega\cubic\meter\per\year}. At the \SI{4}{\percent} BC
Ministry of Forests rate, D26 raises net present value by \SI{36.0}{\millionCAD} and reduces
conditional value at risk at the 20 percent level (CVaR20) by
\SIrange{23.4}{23.7}{\millionCAD}, a net gain of \SIrange{59.4}{59.7}{\millionCAD}. D55 is
approximately break-even (a net \SI{-0.5}{\millionCAD} calibrated and \SI{7.9}{\millionCAD} observed), and D84 is dominated by D55 on net
CVaR. The ranking is rate-dependent: no standard improves net CVaR at \SI{3}{\percent}; all do at
\SIrange{5}{8}{\percent}, and D55 is best at \SI{8}{\percent}. Fire is a scenario, not a forecast;
the observed record and a lower-severity calibrated sample are the two cases we consider. The approach links a regulatory ceiling, enterprise scheduling, composition, and tail risk; at the
BC rate, the AAC need not be traded for risk.

\vspace{1cm}

\paragraph{Keywords} forest policy, deciduous-forward silviculture, harvest scheduling, fire risk,
allowable annual cut, forest economics, British Columbia

\newpage
